\documentclass[11pt]{article}
\usepackage[a4paper,margin=1in]{geometry}
\usepackage{amsmath,amssymb,amsthm,bm,mathtools}
\usepackage{physics}
\usepackage{booktabs}
\usepackage{hyperref}

\title{CV-DV LCHS Formulation Used in \texttt{formed\_qutip.py} and \texttt{formed\_bosonic.py}}
\author{Internal Derivation Note}
\date{\today}

\begin{document}
\maketitle

\section{Problem and Operator Model}
We solve linear ODE systems of the form
\begin{equation}
\frac{d}{dt}\bm{u}(t) = -A\,\bm{u}(t),
\qquad
\bm{u}(T)=e^{-AT}\bm{u}(0).
\end{equation}
In code, $A$ is provided by Pauli decomposition, split into two blocks:
\begin{equation}
L = \sum_j \ell_j P_j,
\qquad
H = \sum_k h_k Q_k,
\qquad
A=L+H.
\end{equation}
(For the 1D heat-equation benchmark, $H=0$ and $L$ is the Laplacian block.)

\paragraph{Heat example used in the files.}
For $2$ DV qubits ($4\times 4$ system),
\begin{equation}
A = \frac{\alpha}{h_{\mathrm{grid}}^2}
\left(2\,II-IX-\frac12XX-\frac12YY\right).
\end{equation}

\section{CV-DV LCHS Representation}
The implemented hybrid map is
\begin{equation}
K(T)=\bra{\phi_r}\,e^{-iT\,\hat{x}\otimes A}\,\ket{\psi_{r,r',\beta}},
\label{eq:lchs-map}
\end{equation}
with
\begin{equation}
\ket{\phi_r}=S(r)\ket{0},
\qquad
\ket{\psi_{r,r',\beta}}=S(r')\sum_{n=0}^{N_c-1}C_n\ket{n}.
\end{equation}
The reference compares $K(T)$ to $e^{-AT}$ up to a global complex scale,
\begin{equation}
\eta^*=\arg\min_\eta\|K-\eta e^{-AT}\|_F,
\qquad
\epsilon_{\mathrm{map}}=\frac{\|K-\eta^*e^{-AT}\|_F}{\|\eta^*e^{-AT}\|_F}.
\end{equation}

\section{Coefficient Formula Used in Code}
\subsection{Kernel branch}
Following the near-optimal branch in arXiv:2312.03916,
\begin{equation}
g_\beta(k)=\frac{\exp\!\left(-(1+ik)^\beta\right)}{C_\beta(1-ik)},
\qquad
C_\beta=2\pi e^{-2^\beta},
\qquad
0<\beta<1.
\end{equation}

\subsection{Corrected $\hbar=1$ expression}
Our QuTiP/bosonic implementation uses
\begin{equation}
\hat{x}=\frac{a+a^\dagger}{\sqrt{2}}
\quad (\hbar=1\text{ convention}).
\end{equation}
The corrected coefficient equation is
\begin{equation}
C_n \propto \int_{-\infty}^{\infty}
H_n\!\left(\frac{k}{\sigma'}\right)
\,g_\beta(k)
\,e^{-\gamma_{\hbar=1}k^2}
\,dk,
\qquad
\gamma_{\hbar=1}=\frac12\left(e^{-2r'}-e^{-2r}\right),
\label{eq:corrected-explicit}
\end{equation}
with $\sigma=e^r$, $\sigma'=e^{r'}$.

Equation \eqref{eq:corrected-explicit} is what \texttt{formed\_qutip.py} labels as the explicit-overlap backend.
The Gauss--Hermite compensated backend is the same formula under
\begin{equation}
k=\sqrt{2}\,\sigma'\xi,
\end{equation}
giving
\begin{equation}
C_n \propto \sum_i w_i
H_n(\sqrt{2}\xi_i)
\,g_\beta(\sqrt{2}\sigma'\xi_i)
\,\exp\!\left(\frac{\sigma'^2}{\sigma^2}\xi_i^2\right).
\label{eq:gh-comp}
\end{equation}

\section{Trotterized Gate Construction in \texttt{formed\_bosonic.py}}
For Pauli terms $A=\sum_m c_mP_m$, the first-order product formula is
\begin{equation}
e^{-iT\hat{x}\otimes A}
\approx
\left(\prod_m e^{-i\Delta t\,c_m\hat{x}\otimes P_m}\right)^{N_t},
\qquad
\Delta t=\frac{T}{N_t}.
\end{equation}
The code compiles each term by basis diagonalization of $P_m$ and parity mapping to one target qubit:
\begin{equation}
e^{-i\Delta t c_m\hat{x}\otimes P_m}
=
B_m^\dagger C_m^\dagger
\,e^{-i\Delta t c_m\hat{x}\otimes Z_t}
\,C_mB_m.
\end{equation}
Because
\begin{equation}
e^{-i\lambda\hat{x}}=D\!\left(-i\lambda/\sqrt{2}\right),
\end{equation}
the primitive gate is implemented with conditional displacement \texttt{cv\_c\_d}.

\section{CV State Preparation Modules and Theory}
\subsection{Injection (simulator baseline)}
\begin{equation}
\ket{\psi_{\mathrm{seed}}}=\sum_n C_n\ket{n}
\end{equation}
is directly loaded by \texttt{cv\_initialize}. This is not a physical pulse sequence, but gives a clean algorithmic baseline.

\subsection{Gate-based SNAP$+$D optimization}
We use the alternating ansatz
\begin{equation}
U_{\mathrm{SNAP+D}}=\prod_{\ell=1}^{L}D(\alpha_\ell)\,\mathrm{SNAP}(\bm{\theta}_\ell),
\end{equation}
optimized for
\begin{equation}
\max_{\{\alpha_\ell,\bm{\theta}_\ell\}}
\left|\bra{\psi_{\mathrm{seed}}}
U_{\mathrm{SNAP+D}}\ket{0}\right|^2.
\end{equation}
This follows the SNAP gate idea of Heeres \emph{et al.}, PRL 115, 137002 (2015), and alternating numerical compilation strategy used in Liu \emph{et al.}, arXiv:2407.10381 (Sec. VI.E).

\subsection{Optimization-free Givens decomposition (Law--Eberly protocol)}
\label{sec:givens}
The Givens method prepares $\ket{\psi_{\mathrm{seed}}}=\sum_{n=0}^{N-1}C_n\ket{n}$
from the vacuum $\ket{0}$ using $N-1$ analytically determined rotation steps,
with \emph{no numerical optimization}.

\paragraph{Givens rotation.}
Define the two-level rotation acting on adjacent Fock levels $\ket{n}$, $\ket{n+1}$:
\begin{equation}
G(n,n{+}1;\theta,\varphi)
=\begin{pmatrix}\cos\theta & -e^{-i\varphi}\sin\theta\\
e^{i\varphi}\sin\theta & \cos\theta\end{pmatrix}
\oplus I_{\perp},
\label{eq:givens-rotation}
\end{equation}
where the $2\times2$ block acts on the $\{\ket{n},\ket{n+1}\}$ subspace
and $I_\perp$ is the identity on all other Fock levels.
Any $N$-component target state can be factored as
\begin{equation}
\ket{\psi_{\mathrm{seed}}}
=G(0,1;\theta_1,\varphi_1)\,G(1,2;\theta_2,\varphi_2)
\cdots G(N{-}2,N{-}1;\theta_{N-1},\varphi_{N-1})\,\ket{0},
\label{eq:givens-decomposition}
\end{equation}
where the angles $\{(\theta_k,\varphi_k)\}_{k=1}^{N-1}$ are computed
recursively from the coefficients $C_n$ by a top-down elimination
(see \texttt{clacod\_heat1d\_stateprep.py}, function \texttt{givens\_decomposition}).
This decomposition is \emph{exact}: the prepared state equals the target
to machine precision, with fidelity $|\braket{\psi_{\mathrm{prep}}}{\psi_{\mathrm{seed}}}|^2=1$.

\paragraph{Law--Eberly physical realization.}
On a coupled qubit--oscillator platform,
each Givens rotation $G(n,n{+}1;\theta,\varphi)$ is realized as one
Law--Eberly step~[5]:
\begin{enumerate}
\item \textbf{Qubit rotation}: $R_y(2\theta_k)\,R_z(\varphi_k)$ on the helper qubit---loads the correct amplitude ratio and phase into the qubit state.
\item \textbf{Number-selective JC $\pi$-swap}: a Jaynes--Cummings interaction
  tuned to the $\ket{n,e}\leftrightarrow\ket{n{+}1,g}$ manifold transfers
  one quantum of excitation from the qubit into the oscillator.
\end{enumerate}
The JC Hamiltonian coupling the qubit to the $n$-th manifold is~[4]
\begin{equation}
H_{\mathrm{JC}} = \hbar\Omega(t)\bigl(\sigma_- a^\dagger + \sigma_+ a\bigr),
\label{eq:jc-hamiltonian}
\end{equation}
which generates transitions $\ket{n,e}\leftrightarrow\ket{n{+}1,g}$ at
Rabi frequency $\Omega_{n}=\Omega_0\sqrt{n+1}$.
Number selectivity is achieved by frequency-tuning the qubit to address
each manifold individually, as demonstrated experimentally by
Hofheinz~\emph{et~al.}~[11] on a superconducting-resonator platform with
states containing up to $\sim10$ photons.

\paragraph{Theorem (Law and Eberly~[5], Eqs.~5--7).}
\emph{Any Fock superposition $\sum_{n=0}^{N-1}C_n\ket{n}$ can be prepared
from $\ket{0,g}$ in exactly $N-1$ qubit-rotation + selective-JC steps.}
The total gate cost is:
\begin{equation}
\text{Resources:}\quad
(N_{\mathrm{active}}-1)\text{ JC pulses}
+\,(N_{\mathrm{active}}-1)\text{ qubit rotations},
\label{eq:givens-resources}
\end{equation}
where $N_{\mathrm{active}}=|\{n:C_n\ne 0\}|$ is the number of occupied Fock levels.

\paragraph{Simulator realization.}
The bosonic-qiskit library provides a \emph{global} JC gate
$U_{\mathrm{JC}}(\theta,\phi)=\exp\bigl(-i\theta(\sigma_- a^\dagger e^{i\phi}+\mathrm{h.c.})\bigr)$,
which couples \emph{all} Fock manifolds simultaneously with
$\sqrt{n+1}$-dependent Rabi frequencies.
A single \texttt{cv\_jc} call therefore cannot implement a
number-selective Givens rotation without crosstalk to other manifolds.
In the current implementation,
the analytically computed state $\ket{\psi_{\mathrm{seed}}}$ is loaded
via \texttt{cv\_initialize} (simulator injection), and the Givens gate
count~\eqref{eq:givens-resources} is reported for physical resource estimation.
On hardware with frequency-selective JC (e.g., superconducting transmon--cavity),
the same angles would translate directly into pulse parameters.

\section{Documented Failure Cases and Lessons}
\subsection{$\hbar=1$ vs $\hbar=2$ inconsistency (historical bug)}
Historical runs mixed a $\hbar=2$ Gaussian rate with a $\hbar=1$
operator/displacement mapping, causing systematic mismatch in $C_n$
and therefore in $K(T)$.
The full diagnosis and conversion rules are consolidated in
Section~\ref{sec:hbar-conventions}; the active implementation now uses
a consistent $\hbar=1$ chain via \eqref{eq:corrected-explicit} and
\eqref{eq:gh-comp}.

\subsection{Why the effect is stronger for heat equation ($H=0$)}
\label{sec:h0-sensitivity}
Let the two position-operator normalizations be related by
\begin{equation}
\hat{x}_{\mathrm{new}}=s\,\hat{x}_{\mathrm{old}},
\qquad
s=\frac{(a+a^\dagger)/\sqrt{2}}{(a+a^\dagger)/2}=\sqrt{2}.
\end{equation}
In the split evolution used for CHO-like analysis,
\begin{equation}
U_{\Delta t}(s)
=e^{-i\Delta t\,s\,\hat{x}\otimes L}\,e^{-i\Delta t\,I\otimes H}.
\end{equation}
First-order expansion gives
\begin{equation}
U_{\Delta t}(s)
=I-i\Delta t\,(I\otimes H+s\,\hat{x}\otimes L)+\mathcal{O}(\Delta t^2),
\end{equation}
so convention mismatch perturbs only the $L$ channel at leading order:
\begin{equation}
\Delta U_{\Delta t}
=-i\Delta t\,(s-1)\,\hat{x}\otimes L+\mathcal{O}(\Delta t^2).
\end{equation}
Therefore:
\begin{itemize}
\item If $H=0$ (heat equation), the entire generator is scaled:
\begin{equation}
U_{\Delta t}(s)=e^{-i\Delta t\,s\,\hat{x}\otimes L}
=U_{\Delta t}^{\mathrm{old}} \text{ with } T\mapsto sT
\;(\text{equiv. }L\mapsto sL).
\end{equation}
Hence the convention change directly shifts diffusion rate and is typically visible.
\item If $\|H\|\gg \|L\|$ (typical CHO setting), the same perturbation competes with an $H$-dominated rotation:
\begin{equation}
\frac{\|\Delta U_{\Delta t}\|}{\|U_{\Delta t}-I\|}
\approx
|s-1|\frac{\|\hat{x}\otimes L\|}{\|I\otimes H+\hat{x}\otimes L\|}
\ll 1.
\end{equation}
Then the mismatch is often seen mostly as amplitude/global-scale drift rather than large waveform-shape change.
\end{itemize}

Equivalent spectral statement for the unsplit map
\begin{equation}
K(T)=\bra{\phi}\,e^{-iT\hat{x}\otimes L}\,\ket{\psi}
\end{equation}
is
\begin{equation}
K(T)=V\,\mathrm{diag}\!\big(m(T\lambda_j)\big)\,V^{-1},
\qquad
m(z)=\bra{\phi}e^{-iz\hat{x}}\ket{\psi},
\end{equation}
for $L=V\mathrm{diag}(\lambda_j)V^{-1}$. Under $\hat{x}\mapsto s\hat{x}$,
\begin{equation}
m(T\lambda_j)\mapsto m(sT\lambda_j),
\end{equation}
so the whole DV spectrum is sampled at rescaled arguments, again showing why $H=0$ is sensitive.

\paragraph{Empirical note in this repository.}
In heat-equation settings ($H=0$), changing only the position-operator
normalization
$\hat{x}=(a+a^\dagger)/\sqrt{2}\leftrightarrow(a+a^\dagger)/2$
while holding other settings fixed typically causes a clear shift in
vector-shape metrics (fidelity and relative vector error). By contrast,
changing only the Gaussian-rate factor
$\gamma_{\hbar=1}\leftrightarrow\gamma_{\hbar=2}$ is often observed
mainly as a global-scale drift in moderate-tail regimes.
The exact magnitude is parameter-dependent and should be reported per run.

\subsection{Failed CV prep by inconsistency aggregation}
A previous internal strategy aggregated inconsistencies over map error, coefficient-backend mismatch, and truncation tail mass. A representative objective is
\begin{equation}
\mathcal{J}(\beta)=
\lambda_1\,\epsilon_{\mathrm{map}}(\beta)
+\lambda_2\,\epsilon_{\mathrm{backend}}(\beta)
+\lambda_3\,\tau_{\mathrm{tail}}(\beta),
\end{equation}
where
\begin{equation}
\epsilon_{\mathrm{backend}}=
\frac{\|\bm{C}^{(\mathrm{exp})}-\eta\bm{C}^{(\mathrm{gh})}\|_2}{\|\eta\bm{C}^{(\mathrm{gh})}\|_2},
\qquad
\tau_{\mathrm{tail}}=\sum_{n=N_c-4}^{N_c-1}|C_n|^2.
\end{equation}
The failure mode was not one single bug, but amplified disagreement when both backend mismatch and truncation pressure were high.

Empirically we often observed better aggregate consistency at lower $\beta$ (within admissible range), because the phase-growth term in
\begin{equation}
(1+ik)^\beta=|1+ik|^\beta e^{i\beta\arg(1+ik)}
\end{equation}
becomes less oscillatory as $\beta$ decreases, which can reduce destructive quadrature cancellation in finite precision.
This is a numerical-stability trend, not a universal theorem; final choice still depends on $r,r',N_c,N_{\mathrm{quad}}$.

\subsection{Current numerical objective in \texttt{formed\_sweep.py}}
The current sweep/optimization workflow treats final-state fidelity and
post-selection probability as \emph{joint} objectives.
For each valid run, we compute
\begin{equation}
F=\text{state fidelity},\qquad
p_{\mathrm{post}}=\|\bm{u}_{\mathrm{DV,post}}\|_2^2.
\end{equation}
Ranking is multi-objective:
\begin{enumerate}
\item Pareto-front rank on $(F,p_{\mathrm{post}})$ (maximize both).
\item Deterministic tie-break by
\begin{equation}
s_{\mathrm{mo}}=\sqrt{F\,p_{\mathrm{post}}}.
\end{equation}
\end{enumerate}
Local continuous refinement minimizes
\begin{equation}
\mathcal{L}_{\mathrm{local}}=1-s_{\mathrm{mo}}.
\end{equation}
In the current default sweep box, parameters are explored over
$r\in[0.5,5]$, $r'\in[0.01,4]$, and $\beta\in[0.3,0.99]$, with the
physical constraint $r'<r$.

\subsection{Number selectivity and the dispersive regime}
\label{sec:jc-selectivity}
The Givens decomposition (Section~\ref{sec:givens}) requires
\emph{number-selective} JC pulses---rotations that act on a single
manifold $\{\ket{n,e},\ket{n{+}1,g}\}$ while leaving all others invariant.
Understanding how this selectivity arises physically---and why it is
absent from gate-level simulators---is central to justifying the
simulation strategy used in this work.

\paragraph{Dispersive Hamiltonian and photon-number-dependent splitting.}
In circuit QED, a transmon qubit coupled to a microwave cavity operates
in the \emph{dispersive regime} when the qubit--cavity detuning is
large compared to the coupling strength,
$|\Delta|=|\omega_q-\omega_c|\gg g$~[13,14].
Adiabatic elimination of the resonant exchange yields the effective
Hamiltonian~[13, Eq.~(44)]:
\begin{equation}
H_{\mathrm{disp}}
=\omega_c\,a^\dagger a
+\frac12\bigl(\omega_q + \chi\bigr)\sigma_z
-\chi\,a^\dagger a\;\sigma_z,
\label{eq:dispersive-H}
\end{equation}
where $\chi=g^2/\Delta$ is the \emph{dispersive shift}.
The qubit transition frequency becomes photon-number-dependent:
\begin{equation}
\omega_q^{(n)}=\omega_q - 2\chi\,n,
\label{eq:number-dependent-freq}
\end{equation}
so consecutive manifolds $\ket{g,n}\leftrightarrow\ket{e,n}$ are
spectroscopically resolved by $2\chi$.
Schuster~\emph{et al.}~[15] experimentally demonstrated this
number-resolved spectroscopy in a superconducting circuit, observing
distinct qubit absorption peaks for each photon number $n=0,1,\ldots,5$.

\paragraph{Selective operations are pulse-level, not gate-level.}
Number selectivity is achieved by driving the qubit with a
\emph{narrow-band pulse} at frequency~$\omega_q^{(n)}$.
If the pulse bandwidth is narrower than the dispersive splitting
($\sigma_\omega\ll 2\chi$), the pulse only rotates the qubit when
exactly~$n$ photons are present, leaving all other manifolds unperturbed.
Hofheinz~\emph{et al.}~[11] used this principle to implement the
Law--Eberly protocol experimentally, synthesizing arbitrary Fock
superpositions with up to $\sim\!10$ photons by alternating
number-selective qubit rotations and resonant JC swaps.

This mechanism is fundamentally a \emph{pulse-level hardware capability}:
\begin{itemize}
\item It requires precise control of the drive frequency and pulse shape.
\item The required pulse duration scales as $T_{\mathrm{pulse}}\gg 1/\chi$
  (typically $\sim\!\mu$s for $\chi/2\pi\sim\text{MHz}$~[15]).
\item It is not expressible as a single entry in a universal gate set;
  rather, it is a composite calibrated waveform~[16].
\end{itemize}

\paragraph{Gate-level simulators cannot reproduce selectivity.}
Gate-level simulators (bosonic-qiskit, hybridlane~[17])
implement the JC interaction as a single global unitary
\begin{equation}
U_{\mathrm{JC}}(\theta,\phi)
=\exp\!\bigl(-i\theta(\sigma_-\,a^\dagger e^{i\phi}+\sigma_+\,a\,e^{-i\phi})\bigr),
\label{eq:global-jc}
\end{equation}
which rotates \emph{all} manifolds simultaneously at the
$\sqrt{n+1}$-dependent Rabi rate.
A pulse calibrated for manifold~$n$
($\theta_n=\pi/(2\sqrt{n+1})$) simultaneously rotates every other
manifold~$m$ by
\begin{equation}
\theta_m^{\mathrm{eff}}=\frac{\pi}{2}\,\frac{\sqrt{m+1}}{\sqrt{n+1}},
\label{eq:jc-crosstalk}
\end{equation}
which is generically incommensurate with $\pi$.
Numerical sanity checks on 8-level targets show that applying
the full Givens sequence via global \texttt{cv\_jc} calls yields
low fidelity with substantial leakage into the qubit-excited sector,
consistent with the crosstalk model above.

To faithfully simulate number-selective operations one would need
a pulse-level simulator (e.g., time-dependent Hamiltonian integration
via QuTiP's \texttt{mesolve} with explicit dispersive-shift parameters),
or a gate set that includes number-selective rotations as primitives.
Neither is available in the current gate-level framework.

\paragraph{Justification for injection in simulation.}
Given this fundamental gap between gate-level simulation and
pulse-level hardware capability, the use of \texttt{cv\_initialize}
(direct Fock-amplitude injection) is the appropriate simulation
strategy for testing the LCHS algorithm.
The reasoning follows a standard separation-of-concerns argument
common in hybrid CV-DV algorithm design~[11,16]:
\begin{enumerate}
\item \textbf{Correctness.}
  Injection loads the exact target state
  $\ket{\psi_{\mathrm{seed}}}=\sum_n C_n\ket{n}$,
  which is precisely the state that a calibrated pulse sequence produces
  on hardware (demonstrated experimentally by Hofheinz~\emph{et al.}~[11]
  with fidelities exceeding 96\% for states up to $n\sim 9$).
\item \textbf{Resource accounting.}
  The Givens decomposition provides the exact pulse count
  \eqref{eq:givens-resources}, which correctly estimates the hardware
  cost without requiring explicit pulse-level simulation.
\item \textbf{Isolation of algorithmic error.}
  The primary goal of the simulation is to validate the LCHS
  approximation quality---Trotter error, post-selection probability,
  and fidelity of $K(T)$ versus~$e^{-AT}$.
  Injection cleanly decouples these algorithmic metrics from
  state-preparation infidelity, which is a device-dependent quantity.
\item \textbf{Precedent.}
  Theoretical studies of bosonic quantum error correction
  (GKP~[18], cat codes~[19]) and hybrid CV-DV algorithms
  routinely assume ideal state preparation in simulation,
  reporting gate counts or pulse counts as separate resource
  metrics.
\end{enumerate}

\section{Practical Recommendations Used by the Two Formed Files}
\begin{itemize}
\item Keep one convention end-to-end ($\hat{x}=(a+a^\dagger)/\sqrt{2}$ in both QuTiP and bosonic mapping).
\item Check explicit-overlap vs GH-comp coefficients before running expensive gate-level circuits.
\item \textbf{CV state preparation strategy} (see Table~\ref{tab:stateprep}):
  \begin{itemize}
  \item \emph{Injection} (\texttt{cv\_initialize}): exact, simulator-only.
    Use as the theoretical baseline against which gate-based methods are compared.
  \item \emph{SNAP$+$D}: gate-compiled, optimization-based.
    Best simulator fidelity ($>99.99\%$ at depth~4).
    Suitable for circuit-level benchmarking.
  \item \emph{Givens}: optimization-free, analytically exact angles.
    Reports the physical JC gate count; uses injection in the simulator
    due to the global-JC selectivity gap (Section~\ref{sec:jc-selectivity}).
    Suitable for hardware resource estimation.
  \end{itemize}
\item Track both map-level error ($\epsilon_{\mathrm{map}}$) and vector-level fidelity after post-selection.
\item In sweep optimization, treat fidelity and post-selection probability as
  joint objectives (Pareto-first ranking with scalar tie-break
  $s_{\mathrm{mo}}=\sqrt{F\,p_{\mathrm{post}}}$).
\item Increase $N_c$ and Fock cutoff together when tail mass is not small.
\end{itemize}

\begin{table}[h]
\centering
\caption{CV state preparation methods implemented in \texttt{formed\_bosonic.py}.}
\label{tab:stateprep}
\begin{tabular}{lcccl}
\toprule
Method & Optimization & Gate-based & Fidelity & Reference \\
\midrule
Injection   & --- & No  & $1.0$ (exact) & --- \\
SNAP$+$D    & L-BFGS-B & Yes & $>0.9999$ (depth~4) & Heeres~[2], Krastanov~[12] \\
Givens      & None     & Yes$^*$ & $1.0$ (analytic) & Law--Eberly~[5], Hofheinz~[11] \\
\bottomrule
\multicolumn{5}{l}{\footnotesize $^*$Gate-compiled on hardware with number-selective JC; simulator uses injection.}
\end{tabular}
\end{table}

\section{Library Convention Notes (with Links)}
\paragraph{QuTiP.}
QuTiP directly documents the single-mode position operator as
\[
x=\frac{1}{\sqrt{2}}(a+a^\dagger).
\]
See:
\begin{itemize}
\item \url{https://qutip.readthedocs.io/en/latest/apidoc/quantumobject.html#qutip.core.operators.position}
\end{itemize}
For CV helper functions, QuTiP also documents the scaling relation
\[
\hbar = \frac{2}{g^2},
\]
with default $g=\sqrt{2}$, giving default $\hbar=1$. See:
\begin{itemize}
\item \url{https://qutip.readthedocs.io/en/stable/_modules/qutip/continuous_variables.html}
\end{itemize}

\paragraph{bosonic-qiskit.}
The bosonic-qiskit gate API does not expose a user-level \texttt{hbar} argument in
\texttt{cv\_d} (or other CV primitives), so convention control is handled by analytic mapping
outside the API. See:
\begin{itemize}
\item \url{https://c2qa.github.io/bosonic-qiskit/c2qa.circuit.CVCircuit.html#c2qa.circuit.CVCircuit.cv_d}
\end{itemize}
Its displacement operator implementation is
\[
D(\alpha)=\exp(\alpha a^\dagger-\alpha^* a),
\]
from source:
\begin{itemize}
\item \url{https://raw.githubusercontent.com/C2QA/bosonic-qiskit/cee0d932e5dd7855dd1eeefc33a488c4c35e9ef0/src/c2qa/operators.py}
\end{itemize}
Therefore, to match a target convention (e.g., $\hat{x}=(a+a^\dagger)/\sqrt{2}$), we must
apply the correct parameter conversion in circuit construction.

\section{Consolidated Oscillator Convention Reference: $\hbar=1$ vs $\hbar=2$}
\label{sec:hbar-conventions}

This section consolidates the historical convention-mismatch diagnosis
and the conversion rules used in the current code path.

There is no universal consensus on the normalization of quadrature operators
in quantum optics.
Weedbrook~\emph{et al.}~[7] identify three common choices $\hbar=1/2,1,2$;
Serafini~[8] devotes Chapter~2 to their systematic comparison.
This section documents which convention the paper uses, which convention
the code uses, and the conversion rules between them.

\subsection{Definition of the two conventions}
\begin{center}
\begin{tabular}{lcc}
\toprule
Quantity & $\hbar=2$ (paper) & $\hbar=1$ (code) \\
\midrule
Position operator
  & $\hat{x}=a+a^\dagger$
  & $\hat{x}=\dfrac{a+a^\dagger}{\sqrt{2}}$ \\[6pt]
Momentum operator
  & $\hat{p}=-i(a-a^\dagger)$
  & $\hat{p}=\dfrac{-i(a-a^\dagger)}{\sqrt{2}}$ \\[6pt]
Commutator
  & $[\hat{x},\hat{p}]=2i$
  & $[\hat{x},\hat{p}]=i$ \\[4pt]
Vacuum wavefunction
  & $\bigl(\frac{2}{\pi}\bigr)^{1/4}e^{-x^2}$
  & $\bigl(\frac{1}{\pi}\bigr)^{1/4}e^{-x^2/2}$ \\[6pt]
Squeezed vacuum ($\sigma=e^r$)
  & $\bigl(\frac{2}{\sigma^2\pi}\bigr)^{1/4}e^{-x^2/\sigma^2}$
  & $\bigl(\frac{1}{\sigma^2\pi}\bigr)^{1/4}e^{-x^2/(2\sigma^2)}$ \\[6pt]
Overlap Gaussian rate
  & $\gamma_{\hbar=2}=e^{-2r'}-e^{-2r}$
  & $\gamma_{\hbar=1}=\frac12\gamma_{\hbar=2}$ \\
\bottomrule
\end{tabular}
\end{center}
The $\hbar=1$ convention matches the standard physics textbook
normalization where position and momentum satisfy $[\hat{x},\hat{p}]=i$
(Braunstein and van Loock~[9], Gerry and Knight~[10]).

\subsection{Where the paper uses $\hbar=2$}
The paper (Das~\emph{et al.}, 2025) never explicitly states a convention;
the $\hbar=2$ choice is implicit in the following equations:
\begin{itemize}
\item \textbf{Paper Eq.\,(6)}: the squeezed vacuum is written as
  $\phi_r(x)=\bigl(\frac{2}{\sigma^2\pi}\bigr)^{1/4}e^{-x^2/\sigma^2}$.
  The factor of~$2$ in the normalization and the exponent $e^{-x^2/\sigma^2}$
  (rather than $e^{-x^2/(2\sigma^2)}$) are the signatures of $\hbar=2$.
\item \textbf{Paper Eq.\,(12)}: the coefficient integral uses
  $\gamma=\sigma'^{-2}-\sigma^{-2}=e^{-2r'}-e^{-2r}$,
  which is $\gamma_{\hbar=2}$.
  In $\hbar=1$, the squeezed-vacuum Gaussian is $e^{-x^2/(2\sigma^2)}$,
  and the overlap rate becomes $\gamma_{\hbar=1}=\gamma_{\hbar=2}/2$.
\item \textbf{Paper Eq.\,(13)}: the prefactor
  $K_n=\sqrt{\sigma/\sigma'}\cdot 1/\sqrt{2^n n!}$
  absorbs the $\hbar=2$ normalization.
  In $\hbar=1$, the equivalent is $1/\sqrt{\sqrt{\pi}\,\sigma'}$
  (see \texttt{fock\_prefactor} in \texttt{formed\_qutip.py}).
\end{itemize}
\emph{The paper is internally self-consistent}---its Tables~1 and~2
validate the analytical coefficients to $10^{-25}$ precision.
The issue arises only when formulas are transferred into $\hbar=1$ code
without adjusting the convention-dependent factors.

\subsection{Conversion rules used in the code}
The production code (\texttt{formed\_qutip.py}, \texttt{formed\_bosonic.py})
applies the following conversions:
\begin{enumerate}
\item \textbf{Gaussian rate}: $\gamma_{\hbar=1}=\frac12(e^{-2r'}-e^{-2r})$
  (\texttt{formed\_qutip.py}, function \texttt{gamma\_hbar1}).
\item \textbf{Position operator}: $\hat{x}=(a+a^\dagger)/\sqrt{2}$
  (\texttt{formed\_qutip.py}, function \texttt{position\_operator\_qutip}).
\item \textbf{Displacement mapping}: $e^{-i\lambda\hat{x}}=D(-i\lambda/\sqrt{2})$
  (\texttt{formed\_bosonic.py}, function \texttt{apply\_pauli\_term\_trotter\_step}).
  In $\hbar=2$ (paper), this would be $D(-i\lambda)$ since $\hat{x}=a+a^\dagger$.
\item \textbf{Kernel constant}: $C_\beta=2\pi e^{-2^\beta}$
  (exponentiation $2^\beta$, convention-independent).
\end{enumerate}

\subsection{Impact summary}
\begin{center}
\begin{tabular}{lccc}
\toprule
Quantity & Paper ($\hbar=2$) & Code ($\hbar=1$) & Factor \\
\midrule
$\gamma$ & $e^{-2r'}-e^{-2r}$ & $\frac12(e^{-2r'}-e^{-2r})$ & $1/2$ \\
Displacement $\alpha$ for $e^{-i\lambda\hat{x}}$ & $-i\lambda$ & $-i\lambda/\sqrt{2}$ & $1/\sqrt{2}$ \\
Vacuum normalization & $(2/\pi)^{1/4}$ & $(1/\pi)^{1/4}$ & $2^{1/4}$ \\
\bottomrule
\end{tabular}
\end{center}

\section*{References}
\begin{enumerate}
\item S. C. Benjamin \emph{et al.}, ``Linear combination of Hamiltonian simulation kernels'' (implementation branch referenced in project as arXiv:2312.03916 for the kernel form).
\item R. Heeres \emph{et al.}, ``Cavity State Manipulation Using Photon-Number Selective Phase Gates,'' \emph{Phys. Rev. Lett.} 115, 137002 (2015).
\item X. Liu \emph{et al.}, ``Gate compilation with alternating SNAP and displacement layers for bosonic modes,'' arXiv:2407.10381.
\item E. T. Jaynes and F. W. Cummings, ``Comparison of quantum and semiclassical radiation theories with application to the beam maser,'' \emph{Proceedings of the IEEE} 51, 89 (1963). \url{https://doi.org/10.1109/PROC.1963.1664}
\item C. K. Law and J. H. Eberly, ``Arbitrary control of a quantum electromagnetic field,'' \emph{Phys. Rev. Lett.} 76, 1055 (1996). \url{https://doi.org/10.1103/PhysRevLett.76.1055}
\item G. Strang, ``On the construction and comparison of difference schemes,'' \emph{SIAM Journal on Numerical Analysis} 5, 506 (1968). \url{https://doi.org/10.1137/0705041}
\item C.~Weedbrook, S.~Pirandola, R.~Garc\'{\i}a-Patr\'{o}n, N.~J.~Cerf, T.~C.~Ralph, J.~H.~Shapiro, and S.~Lloyd, ``Gaussian quantum information,'' \emph{Rev.\ Mod.\ Phys.}\ \textbf{84}, 621 (2012). arXiv:1110.3234.
\item A.~Serafini, \emph{Quantum Continuous Variables: A Primer of Theoretical Methods} (CRC Press, 2017), Ch.~2.
\item S.~L.~Braunstein and P.~van~Loock, ``Quantum information with continuous variables,'' \emph{Rev.\ Mod.\ Phys.}\ \textbf{77}, 513 (2005). arXiv:quant-ph/0410100.
\item C.~C.~Gerry and P.~L.~Knight, \emph{Introductory Quantum Optics} (Cambridge University Press, 2004), Sec.~2.1.
\item M.~Hofheinz, H.~Wang, M.~Ansmann, R.~C.~Bialczak, E.~Lucero, M.~Neeley, A.~D.~O'Connell, D.~Sank, J.~Wenner, J.~M.~Martinis, and A.~N.~Cleland, ``Synthesizing arbitrary quantum states in a superconducting resonator,'' \emph{Nature}\ \textbf{459}, 546 (2009). \url{https://doi.org/10.1038/nature08005}
\item S.~Krastanov, V.~V.~Albert, C.~Shen, C.-L.~Zou, R.~W.~Heeres, B.~Vlastakis, R.~J.~Schoelkopf, and L.~Jiang, ``Universal control of an oscillator with dispersive coupling to a qubit,'' \emph{Phys.\ Rev.\ A}\ \textbf{92}, 040303(R) (2015). \url{https://doi.org/10.1103/PhysRevA.92.040303}
\item A.~Blais, R.-S.~Huang, A.~Wallraff, S.~M.~Girvin, and R.~J.~Schoelkopf, ``Cavity quantum electrodynamics for superconducting electrical circuits: An architecture for quantum computation,'' \emph{Phys.\ Rev.\ A}\ \textbf{69}, 062320 (2004). \url{https://doi.org/10.1103/PhysRevA.69.062320}
\item J.~Koch, T.~M.~Yu, J.~Gambetta, A.~A.~Houck, D.~I.~Schuster, J.~Majer, A.~Blais, M.~H.~Devoret, S.~M.~Girvin, and R.~J.~Schoelkopf, ``Charge-insensitive qubit design derived from the Cooper pair box,'' \emph{Phys.\ Rev.\ A}\ \textbf{76}, 042319 (2007). \url{https://doi.org/10.1103/PhysRevA.76.042319}
\item D.~I.~Schuster, A.~A.~Houck, J.~A.~Schreier, A.~Wallraff, J.~M.~Gambetta, A.~Blais, L.~Frunzio, J.~Majer, B.~Johnson, M.~H.~Devoret, S.~M.~Girvin, and R.~J.~Schoelkopf, ``Resolving photon number states in a superconducting circuit,'' \emph{Nature}\ \textbf{445}, 515 (2007). \url{https://doi.org/10.1038/nature05461}
\item P.~Reinhold, S.~Rosenblum, W.-L.~Ma, L.~Frunzio, L.~Jiang, and R.~J.~Schoelkopf, ``Error-corrected gates on an encoded qubit,'' \emph{Nature Physics}\ \textbf{16}, 822 (2020); see also P.~Eickbusch, V.~Sivak, A.~Z.~Ding, S.~S.~Elder, S.~R.~Jha, J.~Venkatraman, B.~Royer, S.~M.~Girvin, R.~J.~Schoelkopf, and M.~H.~Devoret, ``Fast universal control of an oscillator with weak dispersive coupling to a qubit,'' \emph{Nature Physics}\ \textbf{18}, 1464 (2022). \url{https://doi.org/10.1038/s41567-022-01776-9}
\item B.~Li, S.~Multiverse Computing, ``hybridlane: a PennyLane plugin for hybrid CV-DV simulation,'' Pacific Northwest National Laboratory (2024). \url{https://github.com/pnnl/hybridlane}
\item D.~Gottesman, A.~Kitaev, and J.~Preskill, ``Encoding a qubit in an oscillator,'' \emph{Phys.\ Rev.\ A}\ \textbf{64}, 012310 (2001). \url{https://doi.org/10.1103/PhysRevA.64.012310}
\item M.~Mirrahimi, Z.~Leghtas, V.~V.~Albert, S.~Touzard, R.~J.~Schoelkopf, L.~Jiang, and M.~H.~Devoret, ``Dynamically protected cat-qubits: a new paradigm for universal quantum computation,'' \emph{New J.\ Phys.}\ \textbf{16}, 045014 (2014). \url{https://doi.org/10.1088/1367-2630/16/4/045014}
\end{enumerate}

\end{document}
